\documentclass[11pt,a4paper]{article}
\usepackage[T1]{fontenc}
\usepackage[utf8]{inputenc}
\usepackage{lmodern}
\usepackage[margin=27mm]{geometry}
\usepackage{amsmath,amssymb,amsthm,mathtools}
\usepackage{microtype,booktabs,array,enumitem,xcolor}
\usepackage[colorlinks=true,linkcolor=blue!45!black,citecolor=blue!45!black,urlcolor=blue!50!black,breaklinks=true]{hyperref}
\setlength{\emergencystretch}{3em}
\setlength{\parskip}{0.3em}
\setlist{itemsep=2pt,topsep=4pt}

\newtheorem{theorem}{Theorem}[section]
\newtheorem{lemma}[theorem]{Lemma}
\newtheorem{proposition}[theorem]{Proposition}
\newtheorem{corollary}[theorem]{Corollary}
\newtheorem{maintheorem}{Theorem}
\renewcommand{\themaintheorem}{\Alph{maintheorem}}
\newtheorem{rthm}{Theorem}
\renewcommand{\therthm}{R}
\newtheorem{pprob}{Problem}
\renewcommand{\thepprob}{P}
\theoremstyle{definition}
\newtheorem{definition}[theorem]{Definition}
\newtheorem{example}[theorem]{Example}
\theoremstyle{remark}
\newtheorem{remark}[theorem]{Remark}

\newcommand{\Q}{\mathbb{Q}}
\newcommand{\Z}{\mathbb{Z}}
\newcommand{\N}{\mathbb{N}}
\newcommand{\R}{\mathbb{R}}
\newcommand{\C}{\mathbb{C}}
\newcommand{\F}{\mathbb{F}}
\newcommand{\Qbar}{\overline{\mathbb{Q}}}
\newcommand{\Qt}{\mathbb{Q}_3}
\newcommand{\Qtbar}{\overline{\mathbb{Q}}_3}
\DeclareMathOperator{\Tr}{Tr}
\DeclareMathOperator{\Nm}{N}
\DeclareMathOperator{\ord}{ord}
\DeclareMathOperator{\agcd}{agcd}
\DeclareMathOperator{\Gal}{Gal}
\DeclareMathOperator{\Log}{Log}
\newcommand{\floor}[1]{\left\lfloor #1\right\rfloor}

\hypersetup{pdftitle={Ternary trace dynamics and the transcendence of shifted Mills constants},pdfauthor={Keio}}

\title{Ternary trace dynamics and the transcendence\\ of shifted Mills constants}
\author{Keio}
\date{25 October 2026}

\begin{document}
\maketitle

\begin{abstract}
For a sequence $C=(C_k)_{k\ge1}$ of positive integers with $C_{k+1}\ge 2C_k$, let $\xi(C)$ be the least $A>1$ such that $\floor{A^{C_k}}$ is prime for every $k\ge1$.
Saito proved that $\xi(3^k-1)$ is transcendental if the Riemann hypothesis holds.
He also proved that $\xi(r3^k-1)$ is transcendental unconditionally for every integer $r\ge 4.003\cdot10^{14}$, and remarked that a prime-gap hypothesis would allow this bound to be replaced by~$1$.
We prove unconditionally that $\xi(3^k-1)$ is transcendental.
More generally, $\xi(r3^k-s)$ is transcendental for all odd $r,s$ with $3\nmid s$, $s<3r$ and $13\nmid r/\gcd(r,s)$; in particular $\xi(r3^k-1)$ is transcendental for every odd $r$ not divisible by $13$.
We also prove that $\xi(r3^k+s)$ is transcendental whenever $3\nmid s$, $s\le 3r$, and neither $4s$ nor $13s$ divides $r$. For these schedules the consecutive ratios stay below $3$, so Saito's Type~C criterion does not apply to them, even under the Riemann hypothesis.

The method, which we call ternary trace dynamics, starts from Saito's reduction. If the minimum is algebraic, the reduction produces a cubic Pisot number $\beta$ whose traces along the affine exponent orbit $n_k=C_k/g$ are prime.
It combines four ingredients:
\begin{itemize}
\item a Fermat--Teichm\"uller vanishing of traces modulo~$3$;
\item an affine finite-field return argument, which forces $\Tr(\beta^{n_k})\equiv\pm1\pmod{3^k}$;
\item an identification of the $3$-adic limits of the traces at the fixed point $\lambda=\lim_{3\text{-adic}}n_k$;
\item a Galois transfer of the resulting limit identity to archimedean inequalities.
\end{itemize}
The sign of $\lambda$ governs the method.
For $\lambda>0$ the transfer bounds $\beta$.
For $\lambda<0$ it gives either a lower bound for the small conjugates or a cyclotomic rigidity statement.
For Mills' constant itself $\lambda=0$, and the transfer carries no archimedean information.
For the family $r3^k-1$ the cases our method leaves open are even $r$ and odd multiples of $13$; for both we exhibit explicit infinite families that satisfy all of our $3$-adic constraints.
\end{abstract}

\tableofcontents

%======================================================================
\section{Introduction}\label{sec:intro}

\subsection{Background}
Mills \cite{Mil47} showed that there is a real number $A>1$ such that $\floor{A^{3^k}}$ is prime for every $k\ge1$.
More generally, for a sequence $C=(C_k)_{k\ge1}$ of positive integers with $C_{k+1}\ge2C_k$ let
\[
 \mathcal W(C)=\{A>1:\ \floor{A^{C_k}}\text{ is prime for every }k\ge1\}.
\]
By Saito \cite[Theorem~1.2]{SaiVar}, building on Matom\"aki \cite{Mat10}, this set is non-empty and has a least element. We denote this element by $\xi(C)$, or by $\xi(C_k)$ for an explicit schedule.
Saito proved that Mills' constant $\xi(3^k)$ is irrational \cite{SaiIrr}, and in \cite{SaiVar} he obtained irrationality and transcendence criteria of three types, A, B and~C.
At the critical growth rate $\lim C_{k+1}/C_k=3$ the transcendence question is open in general \cite[Table~1]{SaiVar}.
The transcendence of Mills' constant itself is known only conditionally. It holds under a prime-gap hypothesis with exponent below $\frac12$ \cite[Proposition~5.2]{SaiIrr}, and, by the journal version of \cite{SaiVar}, under the Density Hypothesis.

Two results of \cite{SaiVar} on the shifted family $C_k=r3^k-1$ are relevant here.
\begin{itemize}
\item \cite[Theorem~1.1]{SaiVar}: if the Riemann hypothesis is true, then $\xi(3^k-1)$ is transcendental.
\item \cite[Corollary~2.7]{SaiVar}: $\xi(r3^k-1)$ is transcendental for every integer $r\ge400296054181892\approx4.003\cdot10^{14}$.
Saito adds: ``It is natural to expect that (2.2) with $c=3$ holds for all $x\ge1$. If it were true, then we could replace $4.003\times10^{14}$ with $1$.'' Here (2.2) asserts that $[x,x+3x^{2/3}]$ contains a prime.
The conditional statement is \cite[Corollary~2.8]{SaiVar}; under the Riemann hypothesis it is \cite[Theorem~1.4(C)]{SaiVar}.
\end{itemize}
The bound in Corollary~2.7 comes from an explicit prime-interval theorem, which bounds an algebraic minimum from above.
We use the numbering of arXiv version~3 of \cite{SaiVar}; the journal version is partly renumbered, and there the conditional extension to all $r$ is proved in an appendix.
As far as we can see from the arXiv abstract and the journal version, the schedules treated below appear there either only conditionally ($r3^k-1$ with $r<4.003\cdot10^{14}$) or not at all ($r3^k+s$, and $r3^k-s$ with $s\ge2$). The additional fourth example there concerns $\xi(3^{k-\lfloor(\log k)^{1/2}\rfloor}2^{\lfloor(\log k)^{1/2}\rfloor})$.

\subsection{Main results}
\begin{maintheorem}[Saito's constant, unconditionally]\label{thm:A}
The number $\xi(3^k-1)=\min\{A>1:\floor{A^{3^k-1}}\text{ is prime for all }k\ge1\}$ is transcendental.
More generally, $\xi(r3^k-1)$ is transcendental whenever $r\equiv1,3,9,27,35,81\pmod{104}$.
\end{maintheorem}

\begin{maintheorem}[Negative shifts]\label{thm:B}
Let $r,s$ be odd positive integers with $3\nmid s$ and $s<3r$, and put $d=\gcd(r,s)$.
If $13\nmid r/d$, then $\xi(r3^k-s)$ is transcendental.
In particular, $\xi(r3^k-1)$ is transcendental for every odd $r\ge1$ with $13\nmid r$.
\end{maintheorem}

\begin{maintheorem}[Positive shifts]\label{thm:C}
Let $r,s\ge1$ be integers with $3\nmid s$ and $s\le 3r$.
If $4s\nmid r$ and $13s\nmid r$, then $\xi(r3^k+s)$ is transcendental.
In particular, $\xi(3^k+1)$ and $\xi(3^k+2)$ are transcendental, and so is $\xi(r3^k+s)$ whenever $r/4<s\le 3r$ and $3\nmid s$.
\end{maintheorem}

Theorem~\ref{thm:A} is a special case of Theorem~\ref{thm:B} with an especially short proof, which we give separately in Section~\ref{sec:A}.
For the family $r3^k-1$, Theorem~\ref{thm:B} combined with \cite[Corollary~2.7]{SaiVar} leaves open only even $r$ and odd multiples of~$13$, all below $4.003\cdot10^{14}$.
For example, among $1\le r\le100$ it covers the $46$ odd $r$ other than $13,39,65,91$.
For the open values, our constraints are satisfied by explicit infinite families of cubic Pisot numbers (Section~\ref{sec:residual}), so the method cannot decide them without a new idea.
For the schedules of Theorem~\ref{thm:C} the consecutive ratios $C_{k+1}/C_k$ are less than $3$, so Saito's Type~C theorem does not apply.
His Type~B condition (B6) also fails, for $3$-adic reasons: $v_3(C_k-C_m)=m+v_3(r)$ is bounded in $k$.

\subsection{A new problem: prime traces along affine ternary orbits}
Our approach starts from Saito's reduction (Proposition~\ref{prop:saito}).
If $\xi(C)$ is algebraic, then some power $\beta=\xi^g$ is a cubic Pisot number.
The traces $\Tr(\beta^{C_k/g})$ are the prime values $\floor{\xi^{C_k}}$, and they are approximated by $\beta^{C_k/g}$ from above at a power rate.
This suggests the following problem about Pisot numbers alone; the prime-representing constant no longer appears in it.

\begin{pprob}[Mills shadows on affine ternary orbits]\label{prob:shadow}
Let $\beta$ be a cubic Pisot number, and let $(n_k)$ be an integer sequence with $n_{k+1}=3n_k+\Delta$ for a fixed $\Delta\in\Z$ and $n_k\to\infty$.
When is it possible that both of the following hold for all sufficiently large $k$?
\begin{enumerate}[label=\textup{(S\arabic*)}]
\item\label{S1} $T_k:=\Tr(\beta^{n_k})$ is prime.
\item\label{S2} $0\le\beta^{n_k}-T_k\le T_k^{-17/40}$.
\end{enumerate}
\end{pprob}

The orbit satisfies $n_k=\lambda+R\,3^k$ with $\lambda=-\Delta/2$, so $n_k\to\lambda$ in $\Z_3$.
We call $\lambda$ the \emph{fixed point} of the orbit.
For the schedule $C_k=r3^k+\sigma$ (with $\sigma=-s$ or $\sigma=s$) and $n_k=C_k/g$ we have $\lambda=\sigma/g$.
The case $\lambda=0$ is the exponent $3^k$ of Mills' constant; floors of Pisot powers at these exponents appear in Saito's Problem~1.1 of \cite{SaiInt}.
Our solution of Problem~\ref{prob:shadow} for the relevant orbits is the following.

\begin{rthm}[Ternary trace rigidity]\label{thm:rigidity}
Let $\beta$ be a cubic Pisot number. Let $r\ge1$, $g\in\N$ and $\sigma\in\Z$ with $3\nmid g\sigma$, and suppose $n_k=(r3^k+\sigma)/g\in\N$ for all large $k$.
\begin{enumerate}[label=\textup{(\roman*)}]
\item If \ref{S1} holds, then $T_k\equiv\pm1\pmod{3^{k+v_3(r)}}$ for all large $k$. Moreover, the $3$-adic limits of $T_k$ along residue classes of $k$ satisfy the identity~\eqref{eq:identity} with right-hand side $\pm1$.
\item \textup{(Fermat obstruction.)} If $g\mid2$, $\sigma=-1$, $r\equiv 3^j\pmod{104}$ and $3\nmid\Nm(\beta)$, then \ref{S1} fails.
\item \textup{(Negative fixed point.)} Let $\sigma=-s<0$ with $g\mid 2s$. If \ref{S1} and \ref{S2} hold, then $3\nmid\Nm(\beta)$. If moreover $r$ and $s$ are odd, $g=2$ and $13\nmid r$, then \ref{S1} and \ref{S2} cannot both hold.
\item \textup{(Positive fixed point.)} Let $\sigma=s>0$. If \ref{S1} holds, then $\beta<3^{g/s}$. If moreover $4s\nmid r$ and $13s\nmid r$, then \ref{S1} fails for every admissible $g$. (Every admissible $g$ divides $2s$, since $g\mid C_{k+1}-3C_k=-2s$.)
\end{enumerate}
\end{rthm}

Theorems~\ref{thm:A}--\ref{thm:C} follow from Theorem~\ref{thm:rigidity} and Saito's reduction (Corollary~\ref{cor:shadow}).
The proofs of (ii) and of the first part of (iii) are short (Section~\ref{sec:A}).
The second part of (iii) uses a cyclotomic rigidity argument, plus a check of two explicit numbers when $s=1$ (Section~\ref{sec:B}).
The second part of (iv) uses a finite computation over the $816$ cubic Pisot numbers below~$9$ (Sections~\ref{sec:C} and~\ref{sec:comp}).

\subsection{The method in brief}
Let $\beta_1,\beta_2,\beta_3$ be the roots of the minimal polynomial of $\beta$ in $\Qtbar$, and let $S$ be the set of indices with $|\beta_i|_3=1$.
For $i\in S$ write $\beta_i=\omega_iu_i$, where $\omega_i$ is a root of unity of order prime to $3$ (the Teichm\"uller representative) and $u_i\equiv1$ modulo the maximal ideal.
\begin{enumerate}[label=(\alph*)]
\item \emph{Fermat--Teichm\"uller vanishing.} Every $\omega_i$ has order dividing $104=\operatorname{lcm}(2,8,26)$. Hence $\Tr(\beta^n)\equiv0\pmod3$ whenever $104\mid n$ and $3\nmid\Nm(\beta)$. Along $n_k=(3^k-1)/g$ the exponent is a multiple of $104$ for $k\equiv0\pmod{12}$.
\item \emph{Affine return.} Suppose a prime $p=T_k$ had small $3$-part in the order of $\beta$ modulo $p$. Then the affine orbit would return to $n_k$ modulo that order, so $p$ would divide a later prime $T_{k+h}$. Hence $p\equiv\pm1\pmod{3^{k}}$.
\item \emph{Twisted limits at the fixed point.} Along residue classes of $k$, $T_k$ converges $3$-adically to
\begin{equation}\label{eq:identity}
 L=\sum_{i\in S}\omega_i^{\nu}\,u_i^{\lambda}\qquad(\nu\equiv n_k\bmod104),
\end{equation}
and (b) forces $L=\pm1$.
\item \emph{Galois transfer.} The terms of \eqref{eq:identity} are algebraic numbers whose $g$-th powers are roots of unity times $\beta_i^{\sigma}$. Field isomorphisms transport the identity to $\C$, where it becomes an archimedean statement.
\begin{itemize}
\item For $\lambda>0$ the dominant conjugate gives $\beta^{\lambda}<3$.
\item For $\lambda<0$ the dominant conjugate contributes a term of modulus $<1$. The identity then either bounds the small conjugates from below (combined with \ref{S2} and Baker's theorem), or determines $\beta$ inside a cyclotomic field.
\end{itemize}
\end{enumerate}
The contrast with Mills' constant is visible in (c).
When $\lambda=0$ the principal-unit factors disappear and every term of \eqref{eq:identity} is a root of unity.
The identity then reduces to a congruence condition on $\beta$ modulo $3$, and no archimedean information survives.

\subsection{Relation to previous work}
\begin{itemize}
\item The reduction to cubic Pisot numbers, the floor--trace identity and the fractional-part bound \ref{S2} are due to Saito \cite[Proposition~3.2]{SaiVar}. Saito does not use congruences of traces modulo~$3$, or any $3$-adic argument.
\item The author's unpublished manuscript \cite{K009} treats the exponent $3^r$, i.e.\ the fixed point $\lambda=0$, in connection with composite floors $\floor{\beta^{3^r}}$ (Saito's Problem~1.1 in \cite{SaiInt}). It contains an affine finite-field return and the $3$-adic limit of $\Tr(\beta^{3^r})$ with its Teichm\"uller description. Here we extend the return argument to affine orbits. The twisting by $\omega_i^{\nu}$, the fixed point $\lambda\ne0$, the Fermat obstruction, the Galois transfer and the cyclotomic rigidity are new.
\item The author's unpublished manuscript \cite{K031} gives finite certificates for the four cubic Pisot numbers below~$2$, which lower Saito's threshold to $1.624\cdot10^{14}$. It also contains a late-tail transcendence theorem and the failure of (B6) for $r3^n-1$. It does not treat the case $r=1$.
\end{itemize}
None of the results of \cite{K009,K031,K042} is used in the proofs of Theorems~\ref{thm:A}--\ref{thm:C}, Theorem~\ref{thm:D13} or Proposition~\ref{prop:even}.
To our knowledge, Theorems~\ref{thm:A}--\ref{thm:C} are new.

\paragraph{Organisation.}
\begin{itemize}
\item Section~\ref{sec:saito} recalls Saito's reduction.
\item Section~\ref{sec:dynamics} develops ternary trace dynamics.
\item Sections~\ref{sec:A}, \ref{sec:B} and~\ref{sec:C} prove Theorems~\ref{thm:A}, \ref{thm:B} and~\ref{thm:C}.
\item Section~\ref{sec:residual} describes the remaining cases.
\item Section~\ref{sec:comp} describes the computations.
\end{itemize}

%======================================================================
\section{Saito's reduction}\label{sec:saito}

We write $c_1=C_1$ and $c_{k+1}=C_{k+1}/C_k$.

\begin{proposition}[Saito {\cite[Proposition~3.2]{SaiVar}}]\label{prop:saito}
Let $\theta\in[\frac12,1)$ be such that, for some $d_0,x_0>0$, the interval $[x,x+x^\theta]$ contains at least $d_0x^\theta/\log x$ primes for every $x\ge x_0$.
Let $C_k\in\N$ satisfy $c_1\ge1$, $c_{k+1}\ge2$ for all $k$, and $\limsup c_{k+1}>1/(1-\theta)$.
Assume that $\xi=\xi(C)$ is algebraic and that $\xi^{C_m}\notin\N$ for every $m$.
Fix $0<\varepsilon_0<\limsup c_{k+1}-1/(1-\theta)$ and put $I=\{k:\ c_{k+1}\ge1/(1-\theta)+\varepsilon_0\}$.
Then:
\begin{enumerate}[label=\textup{(\roman*)}]
\item $\{\xi^{C_k}\}\le 2\,c_{k+1}^{-1}\floor{\xi^{C_k}}^{-(1-\theta)c_{k+1}+1}$ for all sufficiently large $k\in I$;
\item there is $g\in\N$ such that $\beta=\xi^g$ is a Pisot number of degree $\ell$ with $2\le\ell\le1+((1-\theta)\limsup c_{k+1}-1)^{-1}$, and $g\mid C_k$ for all sufficiently large $k\in I$;
\item if $\ell=2$, then $\xi^g=(1+\sqrt5)/2$ and $C_k/g$ is an odd prime for all sufficiently large $k\in I$;
\item $\Tr(\beta^{C_k/g})=\floor{\xi^{C_k}}$ for all sufficiently large $k\in I$.
\end{enumerate}
\end{proposition}
By \cite[Theorem~3.1]{SaiVar}, which is derived from the theorem of Baker, Harman and Pintz \cite{BHP01}, $\theta=21/40$ is admissible. Then $1/(1-\theta)=40/19$.
We consider two kinds of schedules:
\begin{equation}\label{eq:schedules}
\begin{aligned}
 &\text{(N)}\quad C_k=r3^k-s &&(r,s\ge1,\ 3\nmid s,\ s<3r),\\
 &\text{(P)}\quad C_k=r3^k+s &&(r,s\ge1,\ 3\nmid s,\ s\le3r).
\end{aligned}
\end{equation}
In both cases $C_k\ge1$ and $c_{k+1}\ge2$ for all $k$, and $c_{k+1}\to3$.
For (N), the condition $c_{k+1}\ge2$ is equivalent to $r3^k+s\ge0$. For (P) it is equivalent to $r3^k\ge s$, i.e.\ to $s\le3r$.
Write $C_k=r3^k+\sigma$ with $\sigma=-s$ or $\sigma=s$. In both cases $3\nmid\sigma$, so $3\nmid C_k$.

\begin{lemma}\label{lem:schedules}
For a schedule of type \textup{(N)} or \textup{(P)}:
\begin{enumerate}[label=\textup{(\alph*)}]
\item for every $m$ there are infinitely many $k$ with $C_m\mid C_k$;
\item for every $g\in\N$, it is not the case that $C_k/g$ is a prime for all sufficiently large $k$;
\item $\gcd(C_k,C_{k+1})$ divides $2s$; hence so does $D:=\lim_m\gcd(C_m,C_{m+1},\dots)$. For $C_k=r3^k-1$ we have $D=2$ if $r$ is odd and $D=1$ if $r$ is even.
\end{enumerate}
\end{lemma}
\begin{proof}
(a) We have $C_k-3^{k-m}C_m=\sigma(1-3^{k-m})$ and $\gcd(3,C_m)=1$, so $C_m\mid C_k$ whenever $\ord_{C_m}(3)\mid k-m$.

(b) Suppose $p_k:=C_k/g$ is prime for all $k\ge k_0$. Take $k\ge k_0$ with $p_k>g$, put $p=p_k$, and let $h=\ord_p(3)$; this is defined since $3\nmid C_k$.
Then $C_{k+h}-C_k=r3^k(3^h-1)\equiv0\pmod p$, so $p\mid C_{k+h}$. As $p\nmid g$, also $p\mid p_{k+h}$. Hence $p_{k+h}=p$, contradicting $C_{k+h}>C_k$.

(c) We have $C_{k+1}-3C_k=-2\sigma$. For $C_k=r3^k-1$, every $C_k$ is even when $r$ is odd, and odd when $r$ is even.
\end{proof}

\begin{lemma}[Scaling]\label{lem:scaling}
For $d\in\N$ we have $\mathcal W(dC)=\{A>1:\ A^d\in\mathcal W(C)\}$, hence $\xi(dC)=\xi(C)^{1/d}$.
In particular, if $d=\gcd(r,s)$, then $\xi(r3^k\pm s)=\xi((r/d)3^k\pm s/d)^{1/d}$, and one of these numbers is transcendental if and only if the other is.
The scaled schedule is again of type \textup{(N)} or \textup{(P)}. In type (N) with $r,s$ odd, the numbers $r/d$ and $s/d$ are again odd.
\end{lemma}
\begin{proof}
We have $\floor{A^{dC_k}}=\floor{(A^d)^{C_k}}$, and $A\mapsto A^d$ is increasing on $(1,\infty)$.
\end{proof}

\begin{corollary}[Mills shadow]\label{cor:shadow}
Let $C$ be of type \textup{(N)} or \textup{(P)}, and suppose that $\xi=\xi(C)$ is algebraic.
Then there are $g\mid D$ and a cubic Pisot number $\beta=\xi^g$ such that $n_k:=C_k/g\in\N$ for all large $k$ and \ref{S1} holds.
For type \textup{(N)}, \ref{S2} holds as well. For type \textup{(P)}, \ref{S2} holds with the exponent $17/40$ replaced by $17/40-\delta$, for any fixed $\delta>0$.
If $C$ is of type \textup{(N)} with $r,s$ odd and $\gcd(r,s)=1$, one may take $g=2$.
\end{corollary}
\begin{proof}
Suppose $\xi^{C_m}=M\in\N$ for some $m$. Then $M\ge2$, and for the infinitely many $k$ of Lemma~\ref{lem:schedules}(a) the number $\floor{\xi^{C_k}}=M^{C_k/C_m}$ is composite, which is impossible.
So Proposition~\ref{prop:saito} applies with $\theta=21/40$, and every large $k$ lies in $I$ because $c_{k+1}\to3>40/19$.
The degree satisfies $\ell\le 1+(57/40-1)^{-1}=1+40/17<4$. The case $\ell=2$ is excluded by (iii) and Lemma~\ref{lem:schedules}(b), so $\ell=3$.
Since $g\mid C_k$ for all large $k$, we have $g\mid D$.
By (iv), $T_k=\floor{\xi^{C_k}}$ is prime, and (i) gives
\[
0\le\{\xi^{C_k}\}=\beta^{n_k}-T_k\le 2c_{k+1}^{-1}T_k^{-(19/40)c_{k+1}+1}.
\]
For (N) we have $c_{k+1}>3$, so the right side is at most $T_k^{-17/40}$. For (P) we only have $c_{k+1}\to3$, which gives the weakened exponent.

For the last claim, note that $C_k$ is even, since $r$ and $s$ are odd.
If $g$ is odd, replace $(g,\beta)$ by $(2g,\beta^2)$. This is allowed for three reasons:
\begin{itemize}
\item $\beta^2$ is a cubic Pisot number, since $\Q(\beta^2)=\Q(\beta)$ by Lemma~\ref{lem:powerfield};
\item $2g\mid C_k$ for all large $k$;
\item the traces do not change: $\Tr_{\Q(\beta^2)/\Q}((\beta^2)^{C_k/2g})=\Tr_{\Q(\beta)/\Q}(\xi^{C_k})$, and $\xi^{C_k}$ itself is unchanged.
\end{itemize}
Now $g=2t$ with $t\mid s$, because $g\mid D\mid2s$. Also $t\mid C_k+s=r3^k$, and since $3\nmid t$ we get $t\mid\gcd(r,s)=1$.
\end{proof}

\begin{lemma}[Power fields]\label{lem:powerfield}
If $\beta$ is a Pisot number and $m\ge1$, then $\Q(\beta^m)=\Q(\beta)$. In particular, if $\beta$ is cubic with non-real conjugates $x,\bar x$, then $x/\bar x$ is not a root of unity.
\end{lemma}
\begin{proof}
An embedding of $\Q(\beta)$ fixing $\beta^m$ sends $\beta$ to a conjugate $\gamma$ with $\gamma^m=\beta^m$. Then $|\gamma|=\beta>1$, so $\gamma=\beta$.
If $(x/\bar x)^m=1$, the three images of $\beta^m$ would not be distinct, contradicting $[\Q(\beta^m):\Q]=3$.
\end{proof}

%======================================================================
\section{Ternary trace dynamics}\label{sec:dynamics}

\subsection{Notation}
Throughout, $\beta$ is a cubic algebraic integer with minimal polynomial
\[
 f(X)=X^3-aX^2+bX-c\in\Z[X],\qquad c=\Nm(\beta)
\]
(not to be confused with the ratios $c_k$ of Section~\ref{sec:saito}).
When $\beta$ is Pisot we also write $\beta$ for its real dominant root, and $x_2,x_3$ for the two other complex roots. Thus $|x_2|,|x_3|<1$ and $|x_2x_3|=|c|/\beta$.
Fix an algebraic closure $\Qtbar$ of $\Qt$ with valuation $v$ normalised by $v(3)=1$, maximal ideal $\mathfrak m$ and residue field $\overline{\F}_3$.
Let $\Qbar\subset\Qtbar$ be the algebraic closure of $\Q$ in it, let $\beta_1,\beta_2,\beta_3\in\Qbar$ be the roots of $f$, and put $S=\{i:\ v(\beta_i)=0\}$.
For $i\in S$, the Teichm\"uller representative $\omega_i$ is the unique root of unity of order prime to $3$ with $\omega_i\equiv\beta_i\pmod{\mathfrak m}$. Then $u_i=\beta_i/\omega_i$ satisfies $v(u_i-1)>0$.
Since $\beta_i\bmod\mathfrak m$ lies in $\F_3$, $\F_9$ or $\F_{27}$, the order of $\omega_i$ divides $2$, $8$ or $26$, hence divides $104$.
The reduction $\bar f=f\bmod3$ has roots $\bar\beta_i$ ($i=1,2,3$), and $\bar\beta_i\ne0$ exactly for $i\in S$.

For an integer sequence $n_k=(r3^k+\sigma)/g$ as in Theorem~\ref{thm:rigidity}, with $3\nmid g\sigma$, put
\[
 \lambda=\sigma/g\in\Z_3,\qquad n_k-\lambda=r3^k/g,\qquad v_3(n_k-\lambda)=k+v_3(r).
\]
The residue $n_k\bmod104$ depends only on $r3^k\bmod104g$, and is therefore purely periodic in $k$.
Let $P$ be a period. For $j\in\Z/P$ let $\nu_j\in\Z/104$ be the common value of $n_k\bmod104$ for $k\equiv j\pmod P$, and let $O_r=\{\nu_j\}$; this set also depends on $\sigma$ and $g$, which are fixed in each application.

\subsection{Principal units}
\begin{lemma}\label{lem:units}
Let $u\in\Qtbar$ with $v(u-1)>0$.
\begin{enumerate}[label=\textup{(\alph*)}]
\item If $v(u-1)>1/2$, then $v(u^3-1)=v(u-1)+1$.
\item There is $j_0=j_0(u)\ge0$ with $v(u^{3^j}-1)\ge j-j_0$ for all $j\ge0$.
\item The map $\Z\to\Qtbar^\times$, $x\mapsto u^x$, extends continuously to $\Z_3$. The extension satisfies $u^{x+y}=u^xu^y$ and $(u^x)^y=u^{xy}$, and $v(u^x-u^y)\ge v_3(x-y)-j_0$ for $x,y\in\Z_3$.
\end{enumerate}
\end{lemma}
\begin{proof}
Write $u=1+t$, so that $u^3-1=3t+3t^2+t^3$.
(a) If $v(t)>1/2$, then $3t$ has valuation $1+v(t)$, and the other two terms have larger valuation.
(b) In general $v(u^3-1)\ge\min(1+v(t),3v(t))$. So the valuation increases at each cubing until it exceeds $1/2$, and from then on it grows by exactly~$1$ per step, by~(a).
(c) For an integer $z$, $u^{3^jz}-1=(u^{3^j}-1)(1+u^{3^j}+\dots)$ has valuation at least $j-j_0$. So $x\mapsto u^x$ is uniformly continuous for the $3$-adic topology on $\Z$, and the identities extend by continuity.
\end{proof}

\subsection{The Fermat--Teichm\"uller obstruction}
\begin{lemma}\label{lem:fermat}
For every $n\ge1$,
\[
 \Tr(\beta^n)\equiv\sum_{i\in S}\bar\beta_i^{\,n}\pmod3,
\]
where the right side, a power sum of the non-zero roots of $\bar f$ counted with multiplicity, lies in $\F_3$.
In particular, if $3\nmid\Nm(\beta)$ and $104\mid n$, then $3\mid\Tr(\beta^n)$.
\end{lemma}
\begin{proof}
We have $\Tr(\beta^n)=\sum_i\beta_i^n$. For $i\notin S$, $\beta_i^n\in\mathfrak m$; and $\mathfrak m\cap\Z=3\Z$.
If $3\nmid c$, then $S=\{1,2,3\}$. For $104\mid n$ each $\bar\beta_i^{\,n}=1$, so the sum is $3=0$ in $\F_3$.
\end{proof}

\subsection{Affine return}
\begin{lemma}[Affine return]\label{lem:return}
Let $\beta$ be a cubic Pisot number and $n_k=(r3^k+\sigma)/g$ with $3\nmid g\sigma$.
If $T_k=\Tr(\beta^{n_k})$ is prime for all sufficiently large $k$, then
\[
 T_k\equiv\pm1\pmod{3^{k+v_3(r)}}\qquad\text{for all sufficiently large }k.
\]
\end{lemma}
\begin{proof}
Since $\beta$ is Pisot, $T_k=\beta^{n_k}+o(1)$ is strictly increasing for large $k$.
Fix a large $k$ such that $p:=T_k$ is prime, $p>3|c|$, and $T_{k+h}$ is prime and greater than $p$ for all $h\ge1$.
In the ring $R_p=\F_p[X]/(f)$ the class $x$ of $X$ is a unit, and $\Tr(\beta^n)\bmod p$ is the trace of multiplication by $x^n$ on $R_p$.
Let $M$ be the multiplicative order of $x$, and write $M=3^{w}M'$ with $3\nmid M'$.
Since $3\nmid gM'$, there is $h\ge1$ with $gM'\mid 3^h-1$. Then
\[
 n_{k+h}-n_k=\frac{r3^k(3^h-1)}{g}
\]
is divisible by $M'$ and by $3^{k+v_3(r)}$.
If $w\le k+v_3(r)$, then $M\mid n_{k+h}-n_k$, so $x^{n_{k+h}}=x^{n_k}$ and $T_{k+h}\equiv T_k\equiv0\pmod p$. As $T_{k+h}$ is a prime larger than $p$, this is impossible.
Therefore $w\ge k+v_3(r)+1$.
The group $R_p^\times$ is the product of a $p$-group and the groups $\F_{p^{f_i}}^\times$, where $f_i\in\{1,2,3\}$ are the degrees of the irreducible factors of $f\bmod p$. So $3^{k+v_3(r)+1}$ divides $p^{f_i}-1$ for some $i$.
If $f_i\in\{1,2\}$, this gives $p\equiv\pm1\pmod{3^{k+v_3(r)+1}}$.
If $f_i=3$, then $3\mid p^3-1$ forces $p\equiv1\pmod3$, so $v_3(p^2+p+1)=1$ and $v_3(p-1)\ge k+v_3(r)$.
\end{proof}

\subsection{Twisted limits at the fixed point}
\begin{lemma}[Interpolation]\label{lem:interp}
Let $\beta$ and $n_k$ be as above; no primality is needed. Put
\[
 L_j=\sum_{i\in S}\omega_i^{\nu_j}\,u_i^{\lambda}\in\Qtbar\qquad(j\in\Z/P),
\]
with $u_i^\lambda$ defined by Lemma~\ref{lem:units}(c). There is $c_0$ such that
$v\bigl(T_k-L_{k\bmod P}\bigr)\ge k+v_3(r)-c_0$ for all large $k$.
\end{lemma}
\begin{proof}
For $i\notin S$ we have $v(\beta_i)\ge1/3$, because $\beta_i$ lies in an extension of $\Qt$ of degree at most $3$. So $v(\beta_i^{n_k})\ge n_k/3$, which exceeds $k+v_3(r)$ for large $k$.
For $i\in S$ and $k\equiv j\pmod P$ we have $\beta_i^{n_k}=\omega_i^{\nu_j}u_i^{n_k}$. By Lemma~\ref{lem:units}(c), $v(u_i^{n_k}-u_i^{\lambda})\ge v_3(n_k-\lambda)-j_0(u_i)=k+v_3(r)-j_0(u_i)$.
\end{proof}

\begin{proposition}[Limit identity]\label{prop:identity}
Let $\beta$ be a cubic Pisot number and $n_k=(r3^k+\sigma)/g$ with $3\nmid g\sigma$, and assume \ref{S1}.
Then $S\ne\emptyset$, and for every $j\in\Z/P$ there is $\varepsilon_j\in\{1,-1\}$ with
\begin{equation}\label{eq:limit}
 \sum_{i\in S}\omega_i^{\nu_j}u_i^{\lambda}=\varepsilon_j .
\end{equation}
Moreover $\varepsilon_j\equiv\sum_{i\in S}\bar\beta_i^{\,\nu_j}\pmod 3$, and this residue is non-zero.
\end{proposition}
\begin{proof}
By Lemmas~\ref{lem:return} and~\ref{lem:interp}, for large $k\equiv j\pmod P$ there are signs $\varepsilon(k)$ with $v(L_j-\varepsilon(k))\ge k+v_3(r)-c_0$.
If two of these signs differed, we would get $v(2)>0$, which is false. So $\varepsilon(k)$ is eventually constant on the class, and letting $k\to\infty$ gives $L_j=\varepsilon_j$.
If $S=\emptyset$, then $L_j=0$, which is impossible.
The congruence follows by reducing modulo $\mathfrak m$, since $u_i^\lambda\equiv1$.
\end{proof}

\subsection{Galois transfer}
For $i\in S$ and $j\in\Z/P$ put
\[
 Y_{i,j}=\omega_i^{\nu_j}u_i^{\lambda}.
\]
By Lemma~\ref{lem:units}(c), $Y_{i,j}^{\,g}=\omega_i^{g\nu_j}u_i^{\sigma}=\zeta_{i,j}\beta_i^{\sigma}$, where $\zeta_{i,j}=\omega_i^{g\nu_j-\sigma}$ is a root of unity. In particular $Y_{i,j}\in\Qbar$.
Let $\Qbar_\C$ be the algebraic closure of $\Q$ in $\C$.

\begin{lemma}[Transfer]\label{lem:transfer}
For every field isomorphism $\psi:\Qbar\to\Qbar_\C$:
\[
 |\psi(Y_{i,j})|=|\psi(\beta_i)|^{\lambda}\qquad(i\in S).
\]
For every $i$ there is such a $\psi$ with $\psi(\beta_i)=\beta$, the dominant root.
For distinct indices $i,i'$, the union of the sets $\{\psi(\beta_{i'}):\psi(\beta_i)=\beta\}$ and $\{\psi(\beta_i):\psi(\beta_{i'})=\beta\}$ is $\{x_2,x_3\}$.
\end{lemma}
\begin{proof}
The first claim follows from $|\psi(\zeta)|=1$ for roots of unity.
Any $\psi$ maps the roots of $f$ in $\Qbar$ bijectively to the roots of $f$ in $\C$.
Composing $\psi$ with elements of $\Gal(\Qbar_\C/\Q)$ realises the action of the Galois group $G\le S_3$ of $f$ on the roots, which is transitive.
For the last claim, fix $\psi_0$ with $\psi_0(\beta_i)=\beta$.
If $G=S_3$, the transposition of $x_2$ and $x_3$ shows that $\psi(\beta_{i'})$ takes both values.
If $G=A_3$, let $\rho$ be the $3$-cycle with $\rho(\beta)=\psi_0(\beta_{i'})$. Then $\rho^{2}\circ\psi_0$ maps $\beta_{i'}$ to $\beta$, and maps $\beta_i$ to $\rho^2(\beta)$, which is the small root different from $\psi_0(\beta_{i'})$.
\end{proof}

\subsection{Approximation from above}
\begin{lemma}[Approximation magnitude]\label{lem:magnitude}
Let $\beta$ be a cubic Pisot number. Suppose there are $A_0>0$, $\kappa\in(0,\frac12)$ and infinitely many $n$ with $|x_2^n+x_3^n|\le A_0\beta^{-\kappa n}$. Then $|\Nm(\beta)|\le\beta^{1-2\kappa}$.
Under \ref{S2} this applies with $\kappa=17/40$, so that $\beta\ge|\Nm(\beta)|^{20/3}$.
\end{lemma}
\begin{proof}
\emph{The totally real case.} Here $|x_2|\ne|x_3|$, since otherwise $x_3=-x_2$ and $\beta=a\in\Z$.
Say $|x_2|>|x_3|$. Then $|x_2^n+x_3^n|\ge|x_2|^n/2$ for large $n$, and $|x_2|^2\ge|x_2x_3|=|\Nm(\beta)|/\beta$.
Hence $|x_2|\le(2A_0)^{1/n}\beta^{-\kappa}$, and letting $n\to\infty$ gives $|\Nm(\beta)|/\beta\le\beta^{-2\kappa}$.

\emph{The complex case.} Write $x_{2,3}=\varrho e^{\pm i\vartheta}$ with $0<\vartheta<\pi$. Then $\varrho^2=|\Nm(\beta)|/\beta$ and $|x_2^n+x_3^n|=2\varrho^n|\cos n\vartheta|$.
Choose $m\in\Z$ with $|n\vartheta-\frac\pi2-m\pi|\le\frac\pi2$. Then
\[
 |\cos n\vartheta|\ge\tfrac2\pi\bigl|n\vartheta-\tfrac\pi2-m\pi\bigr|=\tfrac1\pi\bigl|n\Log\alpha-(2m+1)\Log(-1)\bigr|,
\]
where $\alpha=x_2/x_3$ and $\Log\alpha=2i\vartheta$.
The linear form is non-zero: otherwise $\alpha^{2n}=1$, contrary to Lemma~\ref{lem:powerfield}. Its integer coefficients are $O(n)$.
By Baker's theorem on linear forms in logarithms (for instance in the form of Baker--W\"ustholz \cite{BW93}), it is at least $(2n+2)^{-C'}$, for a constant $C'$ depending on $\alpha$.
Thus $\varrho^n\le A_1n^{C'}\beta^{-\kappa n}$ for a constant $A_1$, and again $\varrho\le\beta^{-\kappa}$.

Under \ref{S2} we have $|x_2^{n_k}+x_3^{n_k}|=\beta^{n_k}-T_k\le T_k^{-17/40}\le2\beta^{-17n_k/40}$ for large $k$.
\end{proof}

\begin{lemma}[Parity]\label{lem:parity}
If $\beta$ is a totally real cubic Pisot number and $n$ is even, then $\beta^n-\Tr(\beta^n)=-(x_2^n+x_3^n)<0$.
Hence \ref{S2} fails along any orbit containing infinitely many even exponents.
\end{lemma}

%======================================================================
\section{Proof of Theorem~\ref{thm:A}}\label{sec:A}

\begin{lemma}\label{lem:208}
If $r\equiv3^{j}\pmod{104}$ for some $j$, then $r\equiv3^{j'}\pmod{208}$ for some $j'$. For $k\equiv-j'\pmod{12}$, both $r3^k-1$ and $(r3^k-1)/2$ are divisible by $104$.
\end{lemma}
\begin{proof}
We have $3^6=729\equiv105=1+104\pmod{208}$. So the subgroup generated by $3$ modulo $208$ has order $12$ and contains both lifts of each power of $3$ modulo $104$.
For $k\equiv-j'\pmod{12}$ we get $r3^k\equiv1\pmod{208}$.
\end{proof}

\begin{lemma}[$3\nmid\Nm(\beta)$ at negative fixed points]\label{lem:3N}
Let $\beta$ be a cubic Pisot number and $n_k=(r3^k-s)/g$, with $s\ge1$, $3\nmid gs$ and $g\le2s$.
If \ref{S1} and \ref{S2} hold, then $3\nmid\Nm(\beta)$.
\end{lemma}
\begin{proof}
Suppose $3\mid\Nm(\beta)$.
Since $\sum_iv(\beta_i)=v(\Nm(\beta))>0$, we have $|S|\le2$. By Proposition~\ref{prop:identity}, $S\ne\emptyset$ and \eqref{eq:limit} holds with $\lambda=-s/g$.
Fix $j$ and write $Y_i=Y_{i,j}$, so that $Y_i^g=\zeta_i\beta_i^{-s}$.
\begin{itemize}
\item If $S=\{i\}$, then $Y_i=\pm1$, so $\beta_i^s=\zeta_iY_i^{-g}$ is a root of unity. This is impossible, because $\psi(\beta_i)=\beta>1$ for a suitable $\psi$.
\item If $S=\{i,i'\}$, then $Y_i+Y_{i'}=\varepsilon_j$. Take $\psi$ with $\psi(\beta_i)=\beta$. By Lemma~\ref{lem:transfer}, $|\psi(Y_i)|=\beta^{-s/g}<1$, hence
\[
 |\psi(\beta_{i'})|^{-s/g}=|\psi(Y_{i'})|=|\varepsilon_j-\psi(Y_i)|<2,\qquad\text{i.e.}\qquad|\psi(\beta_{i'})|>2^{-g/s}\ge\tfrac14 .
\]
Exchanging $i$ and $i'$ and using the last part of Lemma~\ref{lem:transfer}, we get $|x_2|,|x_3|>\frac14$. Hence $|\Nm(\beta)|/\beta=|x_2x_3|>\frac1{16}$, i.e.\ $\beta<16|\Nm(\beta)|$.
Lemma~\ref{lem:magnitude} gives $\beta\ge|\Nm(\beta)|^{20/3}$. Therefore $|\Nm(\beta)|^{17/3}<16$, so $|\Nm(\beta)|=1$, contradicting $3\mid\Nm(\beta)$.\qedhere
\end{itemize}
\end{proof}

\begin{proof}[Proof of Theorem~\ref{thm:A} and of Theorem~\ref{thm:rigidity}\textup{(ii)} and the first part of~\textup{(iii)}]\leavevmode\\
The first part of Theorem~\ref{thm:rigidity}(iii) is Lemma~\ref{lem:3N}, since $g\mid 2s$ implies $g\le2s$.
Now let $r\equiv1,3,9,27,35,81\pmod{104}$. These are the powers of $3$ modulo $104$, and they are all odd.
Suppose that $\xi=\xi(r3^k-1)$ is algebraic.
By Corollary~\ref{cor:shadow} there are $g\in\{1,2\}$ and a cubic Pisot number $\beta=\xi^g$ such that \ref{S1} and \ref{S2} hold for $n_k=(r3^k-1)/g$.
By Lemma~\ref{lem:3N}, $3\nmid\Nm(\beta)$.
By Lemma~\ref{lem:208}, $104\mid n_k$ for infinitely many $k$, and for these $k$ Lemma~\ref{lem:fermat} gives $3\mid T_k$.
Since $T_k\to\infty$, this contradicts \ref{S1}.
This last step uses only \ref{S1}, which proves Theorem~\ref{thm:rigidity}(ii).
\end{proof}

\begin{remark}
The proof uses no prime-gap input beyond the Baker--Harman--Pintz theorem inside Proposition~\ref{prop:saito}, and no bound for the size of $\xi$.
In particular it does not use the explicit prime-interval theorems behind the threshold of \cite[Corollary~2.7]{SaiVar}, and it does not use the Riemann hypothesis.
The only transcendence input is Baker's theorem, which Lemma~\ref{lem:magnitude} uses when $\beta$ is not totally real.
\end{remark}

%======================================================================
\section{Proof of Theorem~\ref{thm:B}}\label{sec:B}

Let $r,s$ be odd with $3\nmid s$ and $s<3r$.
By Lemma~\ref{lem:scaling} we may assume $\gcd(r,s)=1$; the hypothesis of Theorem~\ref{thm:B} then reads $13\nmid r$.
By Corollary~\ref{cor:shadow} we may take $g=2$, so that
\[
 n_k=\frac{r3^k-s}2,\qquad n_{k+1}=3n_k+s,\qquad\lambda=-\tfrac s2,\qquad Y_{i,j}^2=\zeta_{i,j}\beta_i^{-s}.
\]
Since $s$ is odd, the parity of $n_k$ alternates, so $\beta$ is not totally real by Lemma~\ref{lem:parity}.
By Lemma~\ref{lem:3N} (here $g=2\le2s$), $3\nmid\Nm(\beta)$, so $S=\{1,2,3\}$ and $\bar f(0)\neq0$.
We distinguish the factorisation types of $\bar f$ over $\F_3$:
\begin{itemize}
\item[(a)] three roots in $\F_3$, counted with multiplicity;
\item[(b)] $\bar f=(X-e)q(X)$ with $e=\pm1$ and $q$ irreducible of degree $2$;
\item[(c)] $\bar f$ irreducible.
\end{itemize}
By Proposition~\ref{prop:identity}, $p(\nu):=\sum_i\bar\beta_i^{\,\nu}$ is non-zero in $\F_3$ for every $\nu\in O_r$.
The map $n\mapsto3n+s$ on $\Z/8$ has the two cycles
\[
 \mathcal A_s=\{0,\ s,\ 4,\ s+4\}\qquad\text{and}\qquad\mathcal B_s=\{2,\ s+6,\ 6,\ s+2\}\pmod 8,
\]
and $O_r\bmod8$ is one of them.
Note that $\mathcal B_s$ consists of the residues $2$ and $6$ together with two odd residues differing by~$4$.
For $s=1$, we have $O_r\bmod 8=\mathcal B_1$ if and only if $r\equiv5,7,13,15\pmod{16}$; indeed $2n_k+1=r3^k$, and the powers of $3$ modulo $16$ are $1,3,9,11$.

\begin{proposition}[Mod $3$ survivors]\label{prop:mod3types}
If $p(\nu)\ne0$ for all $\nu\in O_r$, then either $\bar f$ is irreducible, or $\bar f=(X-e)(X^2+1)$ with $e=\pm1$ and $O_r\bmod8=\mathcal B_s$.
In the latter case $p(\nu)=-1$ for even $\nu\in O_r$ and $p(\nu)=e$ for odd $\nu\in O_r$.
\end{proposition}
\begin{proof}
(a) For even $\nu$ every root satisfies $\bar\beta_i^{\,\nu}=1$, so $p(\nu)=3=0$. And $O_r$ contains even residues.

(b) Let $\gamma,\gamma^3\in\F_9$ be the roots of $q$. Then $p(\nu)=e^{\nu}+\gamma^\nu+\gamma^{3\nu}$, and $p(0)=3=0$. So $O_r\bmod8=\mathcal A_s$ is impossible, and $O_r\bmod 8=\mathcal B_s$.
\begin{itemize}
\item Suppose $\gamma$ has order $8$, so that $\gamma^4=-1$. For odd $\nu$ the element $\gamma^\nu$ again has order $8$, so its trace $t_\nu=\gamma^\nu+\gamma^{3\nu}$ is $\pm1$. Indeed $(X-\gamma^\nu)(X-\gamma^{3\nu})=X^2-t_\nu X-1$ is irreducible, which excludes $t_\nu=0$.
For the two odd residues $\nu$ and $\nu+4$ of $\mathcal B_s$ we get $p(\nu)=e+t_\nu$ and $p(\nu+4)=e-t_\nu$. One of these vanishes.
\item Suppose $\gamma$ has order $4$. Then $q=X^2+1$ and $\gamma^3=-\gamma$. For $\nu\in\mathcal B_s$ we get $p(\nu)=1-1-1=-1$ if $\nu$ is even (as $\nu\equiv2\pmod4$), and $p(\nu)=e+\gamma^\nu+(-\gamma)^\nu=e$ if $\nu$ is odd.\qedhere
\end{itemize}
\end{proof}

\begin{proposition}[Cyclotomic rigidity]\label{prop:rigidity}
Assume that $13\nmid r$, that \ref{S1} holds along $n_k$, and that $\bar f$ is irreducible.
Then $\Q(\beta)$ embeds into $\Q(\mu_{26})$. In particular $\Q(\beta)$ is abelian, hence totally real.
\end{proposition}
\begin{proof}
Since $\bar f$ is irreducible, $3$ is inert in $K=\Q(\beta)$ and $K\otimes\Qt\cong\Q_{27}$, the unramified cubic extension.
Let $\phi$ be its Frobenius, and label the roots so that $\beta_{m+1}=\phi^m(\beta_1)$ for $m=0,1,2$.
Then $\omega_{m+1}=\phi^m(\omega)$ and $u_{m+1}=\phi^m(u)$, where $\omega=\omega_1\in\mu_{26}$ and $u=u_1$.
Put $y=u^{-s/2}\in\Q_{27}$.
Identity~\eqref{eq:limit} reads
\[
 \sum_{m=0}^{2}\phi^m(\omega^{\nu})\,\phi^m(y)=\varepsilon_\nu\qquad(\nu\in O_r),
\]
a linear system for the vector $(\phi^m(y))_{m}$.
By Dedekind's independence of characters (Artin's lemma), its coefficient matrix $(\phi^m(\omega^\nu))_{\nu,m}$ has rank equal to the $\Qt$-dimension of the span of $\{\omega^\nu:\nu\in O_r\}$.
We claim that this span is all of $\Q_{27}$.

Modulo $13$ the map $n\mapsto3n+s$ has the fixed point $\lambda_{13}\equiv-s/2$, and $n_k-\lambda_{13}\equiv r3^k/2\not\equiv0\pmod{13}$.
So $O_r\bmod13=\{\ell_0+\delta,\ \ell_0+3\delta,\ \ell_0+9\delta\}$, where $\ell_0\in\Z$ represents $\lambda_{13}$ and $\delta\not\equiv0\pmod{13}$.
As $\omega$ has order $13$ or $26$, the elements $\omega^\nu$ ($\nu\in O_r$) are, up to sign and the common factor $\omega^{\ell_0}$, the Frobenius conjugates $\eta,\eta^3,\eta^9$ of a primitive $13$th root of unity $\eta$.
The relations $\sum_ma_m\eta^{3^m}=0$ with $(a_0,a_1,a_2)\in\Qt^3$ form a subspace that is stable under cyclic permutation (apply $\phi$).
Since $X^2+X+1$ is irreducible over $\Qt$, such a subspace is $0$, the line through $(1,1,1)$, the plane $\{a_0+a_1+a_2=0\}$, or everything.
The last two contain $(1,-1,0)$, which would give $\eta=\eta^3$; this is false.
The line would give $\eta+\eta^3+\eta^9=0$, a vanishing $\Q$-linear combination of three distinct primitive $13$th roots of unity. This is impossible, since these roots are linearly independent over $\Q$.
Hence the system has rank $3$, and Cramer's rule gives $y\in\Q(\mu_{26})$.
Then $u^s=y^{-2}$ and $\beta_1^s=\omega^su^s$ lie in $\Q(\mu_{26})$, and $\Q(\beta_1)=\Q(\beta_1^s)$ by Lemma~\ref{lem:powerfield}.
Thus $\Q(\beta_1)\subseteq\Q(\mu_{26})$ is abelian, hence Galois, and a cubic Galois field is totally real.
\end{proof}

\begin{proposition}[The type $(X-e)(X^2+1)$]\label{prop:b0}
Assume that \ref{S1} holds along $n_k$, that $\bar f=(X-e)(X^2+1)$, and that $O_r\bmod8=\mathcal B_s$.
Then $|x_2|,|x_3|>2^{-1/s}$.
If moreover $\beta$ is not totally real and satisfies \ref{S2}, then $|\Nm(\beta)|=1$ and $\beta<4^{1/s}$.
\end{proposition}
\begin{proof}
Here $3=\mathfrak p_1\mathfrak p_2$ in $K$ with residue degrees $1$ and $2$, and $K\otimes\Qt\cong\Qt\times\Q_9$.
Let $\varpi\in\F_9$ with $\varpi^2=-1$.
Label the roots so that $\beta_1\in\Qt$ reduces to $e$, and $\beta_2\in\Q_9$ and $\beta_3=\phi(\beta_2)$ reduce to $\varpi$ and $-\varpi$.
Then $\omega_1=e$, $\omega_2=\iota$ and $\omega_3=-\iota$, with $\iota\in\mu_4$.
Put $y_m=u_m^{-s/2}$, so that $y_m^2=u_m^{-s}=\omega_m^s\beta_m^{-s}$.
The cycle $\mathcal B_s$ contains the even residue $2$ and an odd residue $\nu_0$; put $\iota'=\iota^{\nu_0}=\pm\iota$.
By Propositions~\ref{prop:identity} and~\ref{prop:mod3types},
\[
 y_1-y_2-y_3=-1,\qquad e\,y_1+\iota' y_2-\iota' y_3=e .
\]
Hence
\[
 y_2=\tfrac12\bigl((1+e\iota'^{-1})+(1-e\iota'^{-1})y_1\bigr),\qquad y_3=\tfrac12\bigl((1-e\iota'^{-1})+(1+e\iota'^{-1})y_1\bigr).
\]
Take $\psi$ with $\psi(\beta_1)=\beta$. Then $|\psi(y_m)|=|\psi(\beta_m)|^{-s/2}$, and in particular $|\psi(y_1)|=\beta^{-s/2}<1$.
As $|1\pm e\psi(\iota')^{-1}|=\sqrt2$, we get $|\psi(y_m)|<\sqrt2$ for $m=2,3$.
Since $\{\psi(\beta_2),\psi(\beta_3)\}=\{x_2,x_3\}$, this gives $|x_2|^{-s/2}<\sqrt2$ and $|x_3|^{-s/2}<\sqrt2$, i.e.\ $|x_2|,|x_3|>2^{-1/s}$.
For non-real conjugates this means $|\Nm(\beta)|/\beta=|x_2|^2>4^{-1/s}$, i.e.\ $\beta<4^{1/s}|\Nm(\beta)|$.
Lemma~\ref{lem:magnitude} gives $|\Nm(\beta)|^{20/3}\le\beta<4^{1/s}|\Nm(\beta)|\le4|\Nm(\beta)|$. So $|\Nm(\beta)|=1$ and $\beta<4^{1/s}$.
\end{proof}

\begin{proof}[Proof of Theorem~\ref{thm:B} and of the second part of Theorem~\ref{thm:rigidity}\textup{(iii)}]
Let $r,s$ be odd with $\gcd(r,s)=1$ and $13\nmid r$, and suppose that $\xi(r3^k-s)$ is algebraic.
With $g=2$, $\beta=\xi^2$ is a cubic Pisot number satisfying \ref{S1} and \ref{S2}.
As shown above, $\beta$ is not totally real and $3\nmid\Nm(\beta)$.
By Proposition~\ref{prop:mod3types}, $\bar f$ is either irreducible, or of the form $(X-e)(X^2+1)$ with $O_r\bmod 8=\mathcal B_s$.
The first case contradicts Proposition~\ref{prop:rigidity}.
In the second case, Proposition~\ref{prop:b0} shows that $\beta$ is a cubic Pisot unit with non-real conjugates, $\beta<4^{1/s}$ and $\bar f=(X-e)(X^2+1)$.
\begin{itemize}
\item If $s\ne1$, then $s\ge5$, since $s$ is odd and $3\nmid s$. So $\beta<4^{1/5}<1.32$, which is below the smallest Pisot number $1.3247\ldots$ \cite{Sie44}. Hence there is no candidate, and no computation is needed.
\item If $s=1$, the enumeration of Section~\ref{sec:comp} gives exactly two candidates, the dominant roots of
\[
 X^3-X^2-2X-1\quad(\beta=2.1479\ldots)\qquad\text{and}\qquad X^3-4X^2+X-1\quad(\beta=3.8063\ldots),
\]
both with $\bar f=(X-1)(X^2+1)$.
For these, the twisted limits at $\lambda=-\frac12$ satisfy (Table~\ref{tab:b0})
\[
 L(\nu)\bmod 27\in\{20,13\}\ \text{resp.}\ \{11,19\}\qquad(\nu\equiv2,3\pmod4),
\]
so $L(\nu)\ne\pm1$, contradicting Proposition~\ref{prop:identity}.
\end{itemize}
The same argument, applied directly to $\beta$ with $g=2$, proves the second part of Theorem~\ref{thm:rigidity}(iii); the hypothesis $\gcd(r,s)=1$ is not needed there.
\end{proof}

%======================================================================
\section{Proof of Theorem~\ref{thm:C}}\label{sec:C}

Let $C_k=r3^k+s$ be of type (P), and suppose that $\xi(C)$ is algebraic.
By Corollary~\ref{cor:shadow} there are $g\mid D\mid 2s$ and a cubic Pisot number $\beta=\xi^g$ such that \ref{S1} holds for $n_k=(r3^k+s)/g$.
The fixed point is $\lambda=s/g\ge\frac12$.

\begin{lemma}[Positive fixed point]\label{lem:posbound}
If \ref{S1} holds along $n_k=(r3^k+s)/g$ with $\lambda=s/g>0$, then $\beta<3^{1/\lambda}$.
\end{lemma}
\begin{proof}
By Proposition~\ref{prop:identity}, $S\ne\emptyset$ and $\sum_{i\in S}Y_{i,j}=\pm1$.
Take $i_0\in S$ and $\psi$ with $\psi(\beta_{i_0})=\beta$.
By Lemma~\ref{lem:transfer}, $|\psi(Y_{i_0,j})|=\beta^{\lambda}$, and $|\psi(Y_{i,j})|=|\psi(\beta_i)|^{\lambda}<1$ for $i\ne i_0$.
Hence $\beta^{\lambda}<1+2=3$.
\end{proof}

This proves the first part of Theorem~\ref{thm:rigidity}(iv).
Since $g\mid2s$, $3\nmid s$ and $3\nmid g$, the possible fixed points are $\lambda=s/g\in\frac12\N$, with numerator prime to $3$ in lowest terms. In particular $\lambda\notin\{\frac32,3\}$.
The smallest Pisot number is $1.3247\ldots$, the real root of $X^3-X-1$ \cite{Sie44}.
For $\lambda\ge4$ we have $3^{1/\lambda}<1.32$, so no candidate exists.
For $\frac12\le\lambda\le\frac72$ we have $3^{1/\lambda}\le9$, so $\beta$ belongs to the explicit list of $816$ cubic Pisot numbers below $9$ (Section~\ref{sec:comp}).

\begin{proposition}[Exceptional table]\label{prop:table}
Let $\lambda\in\{\frac12,1,\frac32,2,\frac52,3,\frac72\}$ and let $\beta<3^{1/\lambda}$ be a cubic Pisot number.
The twisted limit $L_\lambda(\nu):=\sum_{i\in S}\omega_i^\nu u_i^\lambda$ equals $\pm1$ for some $\nu\in\Z/104$ only in the cases of Table~\ref{tab:pos}.
In each case the set of such $\nu$ is contained in $\{\nu\equiv\lambda\pmod N\}$, with $N$ as listed.
\end{proposition}

\begin{table}[ht]
\centering
\begin{tabular}{cllc}
\toprule
$\lambda$ & minimal polynomial of $\beta$ & $\beta$ & $N$\\
\midrule
$1$ & $X^3-X^2-1$ & $1.46557$ & $8$\\
$1$ & $X^3-X^2-X-1$ & $1.83929$ & $13$\\
$1$ & $X^3-X^2-2X-1$ & $2.14790$ & $4$\\
$1$ & $X^3-X^2-2X-2$ & $2.26953$ & $13$\\
$1$ & $X^3-X^2-3X-2$ & $2.51155$ & $13$\\
$\tfrac12$ & $X^3-3X^2-X-1$ & $3.38298$ & $13$\\
$\tfrac12$ & $X^3-5X^2-4$ & $5.15077$ & $13$\\
$\tfrac12$ & $X^3-7X^2+5X-4$ & $6.30787$ & $13$\\
$2$ & $X^3-X^2-1$ & $1.46557$ & $8$\\
\bottomrule
\end{tabular}
\caption{All pairs $(\lambda,\beta)$ with $\lambda>0$, $\beta<3^{1/\lambda}$ and $L_\lambda(\nu)=\pm1$ for some $\nu$. There are none for $\lambda\in\{\frac32,\frac52,3,\frac72\}$. The $\lambda=\frac12$ entries are the squares of the second, fourth and fifth $\lambda=1$ entries.}\label{tab:pos}
\end{table}

\begin{proof}[Proof of Theorem~\ref{thm:C} and of Theorem~\ref{thm:rigidity}\textup{(iv)}]
By Lemma~\ref{lem:posbound} and Proposition~\ref{prop:table}, $(\lambda,\beta)$ is an entry of Table~\ref{tab:pos}.
By Proposition~\ref{prop:identity}, every residue $n_k\bmod104$ with $k$ large satisfies $n_k\equiv\lambda\pmod N$.
Here $n_k-\lambda=r3^k/g$; it is an integer for $\lambda\in\{1,2\}$ and an element of $\Z_{(13)}$ for $\lambda=\frac12$. Also $\gcd(3,N)=1$.
\begin{itemize}
\item For $\lambda=1$ we have $g=s$, and $s\mid r$ because $s\mid C_k$. The condition $N\mid r3^k/s$ means $Ns\mid r$ with $N\in\{4,8,13\}$, so $4s\mid r$ or $13s\mid r$.
\item For $\lambda=\frac12$ we have $g=2s$, and the condition $13\mid r3^k/(2s)$ in $\Z_{(13)}$ means $13s\mid r$.
\item For $\lambda=2$ we have $g=s/2$, and the condition $8\mid 2r3^k/s$ means $4s\mid r$.
\end{itemize}
All three are excluded by hypothesis.
For the last statement of Theorem~\ref{thm:C}: if $r<4s$, then neither $4s$ nor $13s$ divides $r$.
\end{proof}

\begin{remark}
The divisibility conditions are exactly those under which some entry of Table~\ref{tab:pos} survives all $3$-adic constraints. So Theorem~\ref{thm:C} is sharp for the method.
For example, for $r=4$, $s=1$ (the schedule $4\cdot3^k+1$) the only surviving candidate is $\beta=\xi$, the dominant root of $X^3-X^2-2X-1$; this is the square of the root of $X^3-X^2-1$.
A single exceptional candidate of this kind might be excluded by a finite-field certificate at another prime. We do not pursue this here.
\end{remark}

%======================================================================
\section{The remaining cases}\label{sec:residual}

\subsection{A family of cubic Pisot units}
\begin{lemma}\label{lem:family}
For every integer $a\ge1$ the polynomial $F_a=X^3-aX^2-X-1$ has exactly one real root $\xi_a$, and $\xi_a\in(a,a+1)$. Its other two roots are non-real with $|x|^2=1/\xi_a$. Hence $\xi_a$ is a cubic Pisot unit, and $\Tr(\xi_a^{-1})=b/c=-1$.
\end{lemma}
\begin{proof}
The coefficient signs $(+,-,-,-)$ have one sign change, so by Descartes' rule $F_a$ has exactly one positive root. It lies in $(a,a+1)$, since $F_a(a)=-a-1<0<a^2+a-1=F_a(a+1)$.
For $x\le-1$ we have $x^3-x\le0$, and for $-1\le x\le0$ we have $x^3-x\le\frac2{3\sqrt3}<1$. Hence $F_a(x)\le(x^3-x)-1<0$ for all $x\le0$.
So $\xi_a$ is the only real root, and the other two roots are complex conjugates with $|x|^2=1/\xi_a<1$. Irreducibility follows since $F_a(\pm1)\ne0$.
\end{proof}

\subsection{Multiples of $13$}
\begin{theorem}\label{thm:D13}
Let $r$ be odd with $13\mid r$.
If $\xi=\xi(r3^k-1)$ is algebraic, then $\xi$ itself is a cubic Pisot number with non-real conjugates.
Moreover $3$ is inert in $\Q(\xi)$, $\Tr(\xi^{-1})=\pm1$ (that is, $b=\pm c$), and $\xi\ge|\Nm(\xi)|^{20/3}$.
\end{theorem}
\begin{proof}
Take $g=2$ and $\beta=\xi^2$, as in Section~\ref{sec:B} with $s=1$.
The arguments there show three things: $\beta$ is not totally real; $3\nmid\Nm(\beta)$; and $\bar f$ is irreducible or of type $(X-e)(X^2+1)$.
Proposition~\ref{prop:b0} and the finite check exclude the second type, since they do not use $13\nmid r$.
So $\bar f$ is irreducible, and we use the notation of the proof of Proposition~\ref{prop:rigidity}.
Now $n_k\equiv6\pmod{13}$ for all $k$, so all rows of the linear system are proportional, and the system reduces to the single identity
\[
 z+\phi(z)+\phi^2(z)=\varepsilon,\qquad z:=\omega^6y,\qquad z^2=\omega^{12}u^{-1}=\varsigma_0\,\beta_1^{-1},
\]
where $\varsigma_0=\omega^{13}=\pm1$.
Put $z_m=\phi^{m-1}(z)$. These are algebraic, and $z_m^2=\varsigma_0/\beta_m$.
Let $E$ be the Galois closure over $\Q$ of the field generated by the $z_m$. It contains all the square roots $\pm z_m$.
For $\tau\in\Gal(E/\Q)$ we have $\tau(z_m)=\pm z_{\pi(m)}$, where $\pi$ is the permutation induced on the $\beta_m$.
Subtracting $\sum_m\tau(z_m)=\varepsilon$ from $\sum_mz_m=\varepsilon$ gives $\sum_{m\in J}z_m=0$, where $J$ is the set of indices where the sign is $-1$.
This is impossible if $|J|=1$ (since $z_m\ne0$), if $|J|=2$ (since $z_a=-z_b$ would give $\beta_a=\beta_b$), and if $|J|=3$ (since $\varepsilon\ne0$).
So $J=\emptyset$, and $\Gal(E/\Q)$ permutes $z_1,z_2,z_3$ exactly as it permutes $\beta_1,\beta_2,\beta_3$.
Hence $z_1\in\Q(\beta_1)$.
In a real embedding of $\Q(\beta_1)$ both $z_1$ and $\beta_1>0$ are real, so $\varsigma_0=z_1^2\beta_1>0$ and $\varsigma_0=1$.
Thus $\delta:=z_1^{-1}\in\Q(\beta_1)$ satisfies $\delta^2=\beta_1$ and $\Tr(\delta^{-1})=\varepsilon$.
In $\C$, $\xi=\pm\delta$ is therefore a cubic Pisot number in $\Q(\beta)$ with $\Tr(\xi^{-1})=\pm1$.
Lemma~\ref{lem:magnitude} applied to $\beta=\xi^2$ gives the last claim.
\end{proof}

\begin{remark}
Let $s$ be odd with $3\nmid s$, $\gcd(r,s)=1$ and $13\mid r$. The same argument then shows: if $\xi(r3^k-s)$ is algebraic, then $\xi$ is a cubic Pisot number with non-real conjugates, $3$ is inert in $\Q(\xi)$, and $\Tr(\xi^{-s})=\pm1$. If $d=\gcd(r,s)>1$ and $13\mid r/d$, Lemma~\ref{lem:scaling} gives the same conclusion for $\xi^{d}=\xi((r/d)3^k-s/d)$, with $s/d$ in place of $s$.
One uses $z=\omega^{\ell_0}u^{-s/2}$ with $2\ell_0+s\equiv0\pmod{13}$, so that $z^2=\pm\beta_1^{-s}$, and then $\xi=\xi^s\cdot(\xi^2)^{-(s-1)/2}\in\Q(\beta)$.
\end{remark}

For $\beta<9$ the computation of Section~\ref{sec:comp} confirms Theorem~\ref{thm:D13}.
The only $\beta<9$ with $L_{-1/2}(\nu)=\pm1$ for some $\nu$ are the numbers $\delta^2$ with
\[
\begin{gathered}
\delta\in\{\,1.32472\ (X^3-X-1),\ \ 1.75488\ (X^3-2X^2+X-1),\ \ 1.83929\ (X^3-X^2-X-1),\\
2.26953\ (X^3-X^2-2X-2),\ \ 2.52138\ (X^3-3X^2+2X-2),\\ 2.24698\ (X^3-2X^2-X+1)\,\}.
\end{gathered}
\]
The good residues are exactly $\nu\equiv6\pmod{13}$.
The last three are excluded: the non-units $2.26953$ and $2.52138$ by Lemma~\ref{lem:magnitude}, and the totally real $2.24698$ by Lemma~\ref{lem:parity}.
The first three are below $2$; the unpublished certificates of \cite{K031}, combined with \cite[Corollary~2.7]{SaiVar}, would exclude them for all $r$, but we do not rely on this.

The conditions of Theorem~\ref{thm:D13} are met by infinitely many cubic Pisot numbers, and for these the $3$-adic constraints hold.
Let $a\ge1$ with $a\equiv1\pmod3$. By Lemma~\ref{lem:family}, $\xi_a$ is a cubic Pisot unit with non-real conjugates and $\Tr(\xi_a^{-1})=-1$. Moreover $F_a\equiv X^3-X^2-X-1$ is irreducible modulo~$3$.
Let $\beta=\xi_a^2$, let $r$ be odd with $13\mid r$, and let $n_k=(r3^k-1)/2$.
Write $\xi_{a,m}=\omega'_mu'_m$ for the $3$-adic conjugates. Since $(r/13)3^k$ is odd, $\omega'^{\,r3^k}_m=(\omega'^{\,13}_m)^{(r/13)3^k}=\omega'^{\,13}_m=:\varsigma_1=\pm1$, the same for all $m$. Hence
\[
 T_k-\varsigma_1\Tr(\xi_a^{-1})=\varsigma_1\sum_m\xi_{a,m}^{-1}\bigl(u'^{\,r3^k}_m-1\bigr).
\]
Since $3$ is unramified, $v(u'_m-1)\ge1$, and therefore $T_k\equiv-\varsigma_1\pmod{3^{k+v_3(r)+1}}$ for every $k$.
In particular, the congruences of Lemma~\ref{lem:return} and the identity~\eqref{eq:limit} both hold, whether or not the $T_k$ are prime.
So the $3$-adic method cannot exclude this family, which is why $13\mid r$ is excluded in Theorem~\ref{thm:B}.

\subsection{Even $r$}
For even $r$, Corollary~\ref{cor:shadow} gives $g=1$, $n_k=r3^k-1$ odd, and $\lambda=-1$.
The case $3\mid\Nm(\xi)$ is excluded by Lemma~\ref{lem:3N}.
But the $3$-adic constraints cannot exclude the following infinite family.

\begin{proposition}\label{prop:even}
Let $r$ be even, $a\equiv2\pmod3$, $a\ge2$, $\beta=\xi_a$ and $n_k=r3^k-1$.
Then $T_k\equiv-1\pmod{3^{k+v_3(r)+1}}$ for every $k\ge1$. The identity~\eqref{eq:limit} holds with $\varepsilon=-1$ for every residue class, whether or not the $T_k$ are prime.
\end{proposition}
\begin{proof}
Modulo $3$ we have $F_a\equiv(X-1)(X+1)^2$. So every $\omega_i=\pm1$ and $\omega_i^{r3^k}=1$, and the conjugates $\beta_i$ lie in extensions of $\Qt$ with ramification index at most~$2$.
Hence $\beta_i^{n_k}=\omega_i^{r3^k}\omega_i^{-1}u_i^{r3^k-1}=\beta_i^{-1}u_i^{r3^k}$, and
\[
 T_k+1=\sum_i\beta_i^{n_k}-\Tr(\beta^{-1})=\sum_i\beta_i^{-1}\bigl(u_i^{r3^k}-1\bigr).
\]
With ramification index at most $2$ we have $v(u_i-1)\ge\frac12$. Expanding $u^3-1=3t+3t^2+t^3$ gives $v(u_i^{3}-1)\ge\frac32$, and Lemma~\ref{lem:units}(a) then gives $v(u_i^{3^k}-1)\ge k+\frac12$.
If $v(w-1)>0$ and $3\nmid m$, then $v(w^m-1)=v(w-1)$. Hence $v(T_k+1)\ge k+v_3(r)+\frac12$, and since $T_k+1\in\Z$, $v_3(T_k+1)\ge k+v_3(r)+1$.
For the identity, $Y_i=\omega_i^{n_k}u_i^{-1}=\omega_i^{r3^k}\beta_i^{-1}=\beta_i^{-1}$, so $\sum_iY_i=\Tr(\beta^{-1})=-1$.
\end{proof}
For instance, $a=2$ gives $\beta=2.54682\ldots$.
To exclude such $\beta$ one must use the primality of $T_k$ at primes other than~$3$, or the sign condition~\ref{S2} more finely. For an individual $\beta$ and $r$ this is sometimes possible with finite-field certificates, but not uniformly.

\subsection{Mills' constant: the fixed point $\lambda=0$}
For $C_k=3^k$ one has $\agcd(C)=\infty$. Saito's reduction gives $\beta=\xi^{3^m}$ \cite{SaiIrr,SaiVar} and $n_k=3^{k-m}$, so $\lambda=0$.
Then $u_i^\lambda=1$, every term of \eqref{eq:limit} is a root of unity, and Lemma~\ref{lem:transfer} gives no information.
The identity reduces to a condition on $\beta$ modulo $3$: the sum of the Teichm\"uller representatives of the non-zero roots of $\bar f$, which Frobenius permutes, must equal $\pm1$. This condition also appears in the unpublished \cite{K009}.
This explains why the method of this paper does not reach Mills' constant: the method is driven by the archimedean size of $u_i^{\lambda}$.
The obstruction is structural. Neither the Fermat obstruction (since $3^{k-m}$ is never divisible by $2$) nor the Galois transfer is available.

%======================================================================
\section{Computations}\label{sec:comp}

All computations use exact integer arithmetic. The programs are in the supplementary archive and are run as \texttt{python3 -I}.
The proofs need them only in two places: Theorem~\ref{thm:B} with $s=1$ (two explicit numbers), and Theorem~\ref{thm:C} (Table~\ref{tab:pos}).
An independent re-implementation with separate code reproduced all of them: the enumeration, Table~\ref{tab:pos} and the lists for $\lambda=-\frac12,-1$ at three different precisions, and Table~\ref{tab:b0} also by a second method using explicitly Hensel-lifted $3$-adic roots.

\paragraph{Enumeration.}
A cubic Pisot number $\beta<B$ has a minimal polynomial with $0\le a\le\floor B+1$, $|b|\le2\floor B+2$ and $1\le|c|\le\floor B$, because $a=\beta+x_2+x_3$, $b=\beta(x_2+x_3)+x_2x_3$ and $c=\beta x_2x_3$.
For every such triple we test irreducibility (absence of rational roots) and the Pisot property by exact Sturm counts:
\begin{itemize}
\item exactly one root lies in $(1,B)$;
\item in the totally real case, exactly two roots lie in $(-1,1)$;
\item in the complex case, the real root exceeds $|c|$.
\end{itemize}
For $B=9$ there are $816$ cubic Pisot numbers: $619$ with non-real conjugates and $197$ totally real.
The two smallest are $1.3247\ldots$ and $1.4656\ldots$, and the next is $1.7548\ldots$.
The non-totally-real units below $4$ with $\bar f=(X-e)(X^2+1)$ are exactly the two of Section~\ref{sec:B}.

\paragraph{Twisted limits.}
Let $\lambda\in\Z_{(3)}$ and $\nu\in\Z/104$.
Take the integer $n$ with $n\equiv\lambda\pmod{3^{30}}$ and $n\equiv\nu\pmod{104}$, increase it by a multiple of $3^{30}\cdot104$, and compute $\Tr(\beta^n)\bmod3^{22}$ by powering the companion matrix.
The result equals $L_\lambda(\nu)\bmod3^{22}$, for two reasons. In an extension of $\Qt$ of ramification index at most $3$, every principal unit satisfies $v(u^{3^j}-1)\ge j$ for $j\ge1$ (Lemma~\ref{lem:units}). And roots of positive valuation contribute $0$ modulo $3^{22}$.
A residue is declared good if the result is $\pm1$.
This can only overestimate the set of $\nu$ with $L_\lambda(\nu)=\pm1$, which is the safe direction for every use above.
The good sets do not change when $(22,30)$ is replaced by $(30,40)$ or by $(45,70)$.
Table~\ref{tab:pos} lists all good cases for $\lambda>0$.
For $\lambda=-\frac12$, the good cases among all $816$ numbers are the six squares listed after Theorem~\ref{thm:D13}, with good set $\{\nu\equiv6\pmod{13}\}$.
For $\lambda=-1$ there are $83$ good numbers, all with $|b|=|c|$.

\begin{table}[ht]
\centering
\begin{tabular}{lcccc}
\toprule
$f$ & $L(2)\bmod243$ & $L(3)\bmod243$ & $L(6)\bmod243$ & $L(7)\bmod243$\\
\midrule
$X^3-X^2-2X-1$ & $101$ & $148$ & $101$ & $148$\\
$X^3-4X^2+X-1$ & $38$ & $208$ & $38$ & $208$\\
\bottomrule
\end{tabular}
\caption{Twisted limits $L_{-1/2}(\nu)$ for the two units of type $(X-1)(X^2+1)$ below $4$. The values depend only on $\nu\bmod 4$. Modulo $27$ they are $20,13$ and $11,19$, so none is $\pm1$.}\label{tab:b0}
\end{table}

\paragraph{Cross-checks.}
\begin{itemize}
\item We verified Lemma~\ref{lem:fermat} on all $600$ candidates with $3\nmid c$, for five exponents divisible by $104$.
\item We recomputed the mod-$3$ survivors of Proposition~\ref{prop:mod3types} by brute force, for $s\in\{1,3,5,7,11,13,17\}$, all odd $r\bmod208$, and all $18$ reductions $\bar f$ with $\bar f(0)\ne0$.
\item For the same data we computed the rank in Proposition~\ref{prop:rigidity} exactly, in $\Z[\zeta_{26}]$ and $\Z[\zeta_8]$. It is $3$ for $13\nmid r$ and drops to $1$ for $13\mid r$, in agreement with Theorem~\ref{thm:D13}.
\item We verified Lemma~\ref{lem:interp} along actual orbits.
\item We checked the congruences for the families in Section~\ref{sec:residual}, also for $\beta>9$.
\end{itemize}
The full outputs are in the archive.

%======================================================================
\begin{thebibliography}{99}

\bibitem{BHP01}
R.~C.~Baker, G.~Harman and J.~Pintz, The difference between consecutive primes, II, \emph{Proc. London Math. Soc. (3)} \textbf{83} (2001), 532--562.

\bibitem{BW93}
A.~Baker and G.~W\"ustholz, Logarithmic forms and group varieties, \emph{J. reine angew. Math.} \textbf{442} (1993), 19--62.

\bibitem{Mat10}
K.~Matom\"aki, Prime-representing functions, \emph{Acta Math. Hungar.} \textbf{128} (2010), 307--314.

\bibitem{Mil47}
W.~H.~Mills, A prime-representing function, \emph{Bull. Amer. Math. Soc.} \textbf{53} (1947), 604.

\bibitem{SaiIrr}
K.~Saito, Mills' constant is irrational, \emph{Mathematika} \textbf{71} (2025), e70027; arXiv:2404.19461.

\bibitem{SaiVar}
K.~Saito, Transcendency of variants of Mills' constant, arXiv:2508.16068v3 (2025). Journal version: \emph{Ramanujan J.} \textbf{70} (2026), article~69, with partly different numbering. Theorem numbers in this paper refer to arXiv version~3.

\bibitem{SaiInt}
K.~Saito, Intervals without primes near an iterated linear recurrence sequence, \emph{Int. J. Number Theory} (2026); arXiv:2504.14968.

\bibitem{Sie44}
C.~L.~Siegel, Algebraic integers whose conjugates lie in the unit circle, \emph{Duke Math. J.} \textbf{11} (1944), 597--602.

\bibitem{K009}
Keio, Composite terms in sparse cubic Pisot floors: alternating shifts, $3$-adic limits, and projective return, unpublished research manuscript (September 2026), item [009] of the Verifiable Mathematical Reasoning Corpus, \url{https://github.com/QEEDDD/Verifiable-Mathematical-Reasoning-Corpus}.

\bibitem{K031}
Keio, Prime-representing constants and related structures: encodings, arithmetic, geometry, and modular obstructions, unpublished research manuscript (2026), item [031] of the same corpus.

\bibitem{K042}
Keio, Cubic Pisot floors: composite windows, sparse exceptions, and finite-sieve barriers, unpublished research manuscript (October 2026), item [042] of the same corpus.

\end{thebibliography}

\end{document}
